\section{Agent 可读性：让代码库对 Agent ``可见''}

\subsection{以 Agent 可读性为第一优先}

因为仓库完全由 Agent 生成，它首先为 \textbf{Codex 的可读性}做优化——而不是人类的阅读习惯。这类似于团队为新入职工程师优化代码可导航性，只是对象换成了 Agent。

\begin{importantbox}{Agent 可读性的核心原则}
从 Agent 的视角看，\textbf{任何它在运行时无法在上下文中访问的东西，实际上都不存在}。存放在 Google Docs 里的讨论、Slack 聊天记录、或人们脑中的决策——这些对 Agent 都是不可见的。只有仓库本地的、版本化的产出物（代码、Markdown、Schema、可执行计划）才是 Agent 能看到的。
\end{importantbox}

这意味着团队需要不断将更多上下文\textbf{推入仓库}中。那次在 Slack 上对齐团队架构方向的讨论？如果 Agent 发现不了，它就像一个三个月后加入的新人——对此一无所知。

\subsection{提升应用可读性的具体措施}

随着代码产出量的增加，瓶颈转移到了人类的 QA 能力上。团队通过以下方式让更多能力直接对 Agent 可见：

\begin{table}[H]
\centering
\small
\begin{tabular}{p{4cm}p{9cm}}
\toprule
\textbf{措施} & \textbf{效果} \\
\midrule
每个 Git worktree 可独立启动应用 & Codex 可以为每个变更启动和驱动一个独立实例 \\
接入 Chrome DevTools Protocol & Agent 可以操作 DOM 快照、截图和页面导航 \\
本地可观测性栈 & 每个 worktree 有独立的日志、指标和追踪，任务完成后销毁 \\
LogQL / PromQL 查询能力 & Agent 可以直接查询日志和指标 \\
\bottomrule
\end{tabular}
\caption{提升 Agent 可读性的关键措施}
\end{table}

这些能力使得以下类型的 Prompt 成为可能：
\begin{itemize}[nosep]
  \item ``确保服务启动在 800ms 内完成''
  \item ``这四个关键用户旅程中没有任何 span 超过两秒''
\end{itemize}

Agent 可以直接复现 Bug、验证修复、并对 UI 行为进行推理——这些曾经需要人类手动完成的 QA 工作现在被自动化了。

\subsection{技术选型偏好}

\begin{knowledgebox}{``无聊技术''对 Agent 更友好}
团队发现，通常被描述为``无聊的''（boring）技术往往更容易被 Agent 建模，原因包括：
\begin{itemize}[nosep]
  \item \textbf{可组合性}（composability）：简单接口易于推理
  \item \textbf{API 稳定性}：不会因为版本变化而打破 Agent 的既有认知
  \item \textbf{训练集覆盖率}：经典技术在 LLM 训练数据中有更充分的表示
\end{itemize}
在某些情况下，让 Agent 重新实现功能子集比使用不透明的第三方库更划算。例如，团队没有使用通用的 \texttt{p-limit} 式并发控制包，而是自行实现了一个并发映射辅助函数，它与 OpenTelemetry 监控紧密集成、100\% 测试覆盖、且行为完全符合运行时预期。
\end{knowledgebox}

\subsection{本章小结}

Agent 可读性是 Agent-First 开发的基础设施层。核心原则是：Agent 在运行时看不到的东西等于不存在。团队需要持续将讨论、决策、知识推入版本化的仓库产出物中，并通过工具链（DevTools、可观测性栈等）让 Agent 能直接与应用交互。技术选型上，优先选择``无聊但稳定''的技术。

\newpage

